\documentclass[10pt,a4paper]{amsart}
\usepackage[margin=19mm]{geometry}
\usepackage{amsmath,amssymb,mathtools}
\usepackage[scr=rsfs,cal=euler]{mathalfa}
\usepackage{xfrac}
\usepackage[unicode,hidelinks,bookmarks=false]{hyperref}
\newcommand{\ie}{i.e.\ }
\newcommand{\cf}{cf.\ }
\newcommand{\resp}{respectively\ }
\newcommand{\Subsection}{\subsection}
\providecommand{\nobreakdash}{\nobreak\hbox{-}}
\setlength{\emergencystretch}{2em}
\multlinegap0pt
\newcommand{\ee}{\mathrm{e}}
\usepackage{amssymb}
\usepackage{mathtools}
\usepackage[shortlabels]{enumitem}
\newtheorem{proposition}{Proposition}
\newcommand{\C}{\mathbb{C}}
\newcommand{\N}{\mathbb{Z}_{\ge 0}}
\newcommand{\Z}{\mathbb{Z}}
\title{The holonomic triangle: from a symmetry between $\ee$ and $\pi$ to additive Gamma functions}
\author{Benoit Cloitre}
\date{Source: arXiv:2601.04242v1 (CC BY 4.0)}
\begin{document}
\maketitle

\noindent\emph{Source and licence.} The holonomic triangle: from a symmetry between $\ee$ and $\pi$ to additive Gamma functions. Version arXiv:2601.04242v1 (CC BY 4.0), distributed by arXiv under the Creative Commons Attribution 4.0 International licence, http://creativecommons.org/licenses/by/4.0/, as recorded in the arXiv licence metadata for this exact version and shown on the abstract page https://arxiv.org/abs/2601.04242v1. Full licence notice: \texttt{licenses/arxiv-2601.04242-CC-BY-4.0.txt}.\par\smallskip

\noindent\textit{Editorial scope.} Counted unit: the article's proposition prop:g-explicit (source0.tex lines 625--630). The article's complete original proof of it is delivered with it (source0.tex lines 632--685, one \texttt{proof} environment of 54 lines). No mathematics is written, repaired or re-derived: the statement and the proof are the article's own lines, in the article's own notation, and no retained line is substituted or re-flowed.  Retained uncounted as the source-local context the proof needs: lines 617--619; lines 689--689.  Nothing was omitted from any retained span, so the package carries no omission record.  No retained line contains a citation, so the wrapper carries no bibliography and no bibliography entry.  No retained line contains a figure environment, an image-inclusion command or drawing code, so no image is delivered and the package declares no asset dependency.  Editorial formatting only, in addition: an AMS \texttt{amsart} preamble that restates the article's own declarations and macros used by the retained lines, namely the environments \texttt{proposition} on separate counters, together with the standard packages those lines need.  The retained environments print in this document as Proposition 1 in order of appearance, whereas the article prints the same items with the numbering of its own full document.  The complete native source of the article as distributed in the arXiv e-print for this exact version is bundled unchanged as \texttt{sources/192284/source0.tex}.
\begin{equation}\label{eq:g-AFE}
g(z+2) = g(z) - \frac{g(z+1)}{z+1}.
\end{equation}

\section{The AGF prototypes: $\Gamma$, $f$, and $g$}\label{sec:gamma-prototype}

\begin{proposition}\label{prop:g-explicit}
The connection constant $g(z)$ has simple poles at the negative integers $z \in \{-1,-2,-3,\ldots\}$ and is holomorphic at $z=0$. It is given by:
\begin{equation}\label{eq:g-explicit}
g(z) = \sqrt{2}\left[\frac{\Gamma\left(\frac{z}{2}+1\right)}{\Gamma\left(\frac{z+1}{2}\right)} - \frac{\Gamma\left(\frac{z+1}{2}\right)}{\Gamma\left(\frac{z}{2}\right)}\right].
\end{equation}
\end{proposition}

\begin{proof}
We prove this identity by showing that the proposed expression satisfies the functional equation \eqref{eq:g-AFE} and the initial conditions. Since both the integral definition and the proposed formula satisfy the same order-2 AFE with exponential type $<\pi$, agreement at $z=0$ and $z=1$ (verified below) propagates to all of $\Z$ via the recurrence. By Carlson's theorem, the difference of two functions with exponential type $<\pi$ that vanishes on $\N$ must vanish identically on $\C$.

We first verify that the proposed formula \eqref{eq:g-explicit} has exponential type $<\pi$ in vertical strips, as required for uniqueness. The uniform version of Stirling's formula shows that for Gamma ratios of the form $\Gamma(z/2+\alpha)/\Gamma(z/2+\beta)$, the dominant exponential factors $\exp(-\pi|\Im z|/4)$ cancel exactly, leaving polynomial growth (specifically, decay $O(|z|^{-1/2})$) in $|z|$. Since polynomial growth implies exponential type $0 < \pi$, the growth condition required for uniqueness is satisfied. (A detailed derivation is given in Section~\ref{sec:gamma-prototype}.)

Define the auxiliary function:
$$A(z) = \frac{\Gamma\left(\frac{z}{2}+1\right)}{\Gamma\left(\frac{z+1}{2}\right)}.$$

We have $A(z-1) = \frac{\Gamma\left(\frac{z-1}{2}+1\right)}{\Gamma\left(\frac{z-1+1}{2}\right)} = \frac{\Gamma\left(\frac{z+1}{2}\right)}{\Gamma\left(\frac{z}{2}\right)}$.
Thus, the proposed formula is $g(z) = \sqrt{2}(A(z) - A(z-1))$.

We derive relations between shifted values of $A(z)$ using $\Gamma(x+1)=x\Gamma(x)$:
\begin{align*}
A(z)A(z-1) &= \frac{\Gamma\left(\frac{z}{2}+1\right)}{\Gamma\left(\frac{z+1}{2}\right)} \cdot \frac{\Gamma\left(\frac{z+1}{2}\right)}{\Gamma\left(\frac{z}{2}\right)} = \frac{\frac{z}{2}\Gamma\left(\frac{z}{2}\right)}{\Gamma\left(\frac{z}{2}\right)} = \frac{z}{2}.
\end{align*}

This implies $A(z+1)A(z) = \frac{z+1}{2}$, so $A(z+1) = \frac{z+1}{2A(z)}$.
Also, $A(z+2) = \frac{z+2}{2A(z+1)} = \frac{z+2}{2} \frac{2A(z)}{z+1} = \frac{z+2}{z+1}A(z)$.

We now verify the functional equation $g(z+2) = g(z) - \frac{g(z+1)}{z+1}$. Dividing by $\sqrt{2}$, this becomes:
$$A(z+2)-A(z+1) = A(z)-A(z-1) - \frac{A(z+1)-A(z)}{z+1}.$$

Rearranging, we need to verify:
$$A(z+2)+A(z-1) = A(z+1)\frac{z}{z+1} + A(z)\frac{z+2}{z+1}.$$

We compute both sides using the relations derived above. Left-hand side:
$$A(z+2)+A(z-1) = \frac{z+2}{z+1}A(z) + \frac{z}{2A(z)}.$$

Right-hand side:
$$A(z+1)\frac{z}{z+1} + A(z)\frac{z+2}{z+1} = \frac{z+1}{2A(z)}\frac{z}{z+1} + \frac{z+2}{z+1}A(z) = \frac{z}{2A(z)} + \frac{z+2}{z+1}A(z).$$

Since the left-hand side equals the right-hand side, the functional equation is satisfied. Finally, we verify the initial values match those derived from the integral $L_m$. We have $L_1 = \pi/2-1$, hence 
$$g(1) = \sqrt{2/\pi}(\pi/2-1) = \frac{\pi-2}{\sqrt{2\pi}}.$$

We check the formula \eqref{eq:g-explicit} at $z=1$. 
$$A(1) = \frac{\Gamma(3/2)}{\Gamma(1)} = \sqrt{\pi}/2$$ 
and 
$$A(0) = \frac{\Gamma(1)}{\Gamma(1/2)} = 1/\sqrt{\pi}.$$

Then 
$$ \sqrt{2}(A(1)-A(0)) = \sqrt{2}(\sqrt{\pi}/2 - 1/\sqrt{\pi}) = \sqrt{2}\frac{\pi-2}{2\sqrt{\pi}} = \frac{\pi-2}{\sqrt{2\pi}}=g(1).$$
This matches.  We also verify a second initial value to ensure complete uniqueness. For $z=0$: $L_0 = 1$, hence 
$$g(0) = \sqrt{2/\pi} \cdot 1 = \sqrt{2/\pi}.$$

Using the explicit formula \eqref{eq:g-explicit}, we have
$$g(0) = \sqrt{2}(A(0)-A(-1)).$$
We compute
$$A(0) = \frac{\Gamma(1)}{\Gamma(1/2)} = \frac{1}{\sqrt{\pi}}$$
and
$$A(-1) = \lim_{z \to -1} A(z) = \lim_{z \to -1} \frac{\Gamma(z/2+1)}{\Gamma((z+1)/2)} = 0.$$
Thus
$$g(0) = \sqrt{2}(1/\sqrt{\pi} - 0) = \sqrt{2/\pi},$$
confirming the expected value. Since the expression \eqref{eq:g-explicit} satisfies the same linear functional equation as the connection constant $g(z)$, matches the initial conditions, and exhibits the required growth properties, the identity is proven by uniqueness.
\end{proof}

\end{document}
